\documentclass[12pt, a4paper]{article}

% --- PAQUETES ---
\usepackage[utf8]{inputenc}
\usepackage[spanish]{babel}
\usepackage{geometry}
\usepackage{amsmath} % Para fórmulas matemáticas
\usepackage{amsfonts} % Para fuentes matemáticas
\usepackage{amssymb} % Para símbolos matemáticos
\usepackage{graphicx} % Para incluir imágenes (si es necesario)
\usepackage{enumitem} % Para personalizar listas
\usepackage{hyperref} % Para enlaces

% --- CONFIGURACIÓN DE PÁGINA Y TÍTULOS ---
\geometry{a4paper, margin=1in}
\hypersetup{
    colorlinks=true,
    linkcolor=blue,
    filecolor=magenta,      
    urlcolor=cyan,
}

\begin{document}

\title{\bfseries Problemas de Práctica: Semana 1 \\ \large Medición, Resolución de Problemas y Cinemática}
\author{Basado en "Física, 6ta Edición" de Wilson, Buffa y Lou}
\date{}
\maketitle
\thispagestyle{empty}
\vspace{-1cm}

\section*{Instrucciones}
Resuelva los siguientes problemas, que han sido seleccionados de los capítulos 1 y 2. Están ordenados por dificultad creciente. Intente resolverlos aplicando los conceptos y las estrategias de resolución de problemas vistos en la semana.

\hrule
\vspace{1cm}

\begin{enumerate}[leftmargin=*, label=\textbf{Problema \arabic*.}, itemsep=2em]

    \item \textbf{(Dificultad: Muy Fácil) Conversión de unidades.} \\
    La rapidez de la luz en el vacío es de aproximadamente $3.00 \times 10^8$ m/s. Convierta esta rapidez a millas por hora (mi/h). \\
    \textit{(Sugerencia: 1 mi = 1609 m)}

    \item \textbf{(Dificultad: Fácil) Cifras significativas.} \\
    Realice las siguientes operaciones y exprese las respuestas con el número correcto de cifras significativas:
    \begin{enumerate}[label=\alph*)]
        \item $1.24 \times 8.2 = $
        \item $6.78 / 1.2 = $
        \item $1.99 + 2.0 = $
    \end{enumerate}

    \item \textbf{(Dificultad: Fácil) Análisis dimensional.} \\
    Demuestre que la ecuación $x = x_0 + v_0 t + \frac{1}{2}at^2$ es dimensionalmente correcta, donde $x$ y $x_0$ son longitudes, $v_0$ es velocidad, $a$ es aceleración y $t$ es tiempo.

    \item \textbf{(Dificultad: Media-Baja) Rapidez promedio.} \\
    Un estudiante conduce 105 km de su casa a la universidad con una rapidez promedio de 95 km/h. En el viaje de regreso, recorre la misma distancia, pero su rapidez promedio es de 105 km/h. ¿Cuál es la rapidez promedio del estudiante para el viaje redondo completo?

    \item \textbf{(Dificultad: Media) Aceleración constante.} \\
    Un automóvil que viaja a 95 km/h frena a una rapidez constante hasta detenerse en 40 m. ¿Cuál es su aceleración en m/s$^2$?

    \item \textbf{(Dificultad: Media) Caída libre desde el reposo.} \\
    La Torre Inclinada de Pisa mide 54.5 m de altura. Si se deja caer una bola de plomo desde lo alto de la torre, ¿cuánto tiempo tardará en llegar al suelo?

    \item \textbf{(Dificultad: Media-Alta) Movimiento en dos etapas.} \\
    Un automóvil que viaja a 60.0 km/h frena hasta detenerse con una desaceleración uniforme de 4.00 m/s$^2$.
    \begin{enumerate}[label=\alph*)]
        \item ¿Cuánto tiempo tarda en detenerse?
        \item ¿Qué distancia recorre en ese tiempo?
    \end{enumerate}

    \item \textbf{(Dificultad: Alta) Caída libre con velocidad inicial.} \\
    Se lanza una pelota hacia arriba con una rapidez inicial de 20.0 m/s.
    \begin{enumerate}[label=\alph*)]
        \item ¿Cuánto tiempo está en el aire? (Es decir, en su trayecto de ida y vuelta).
        \item ¿Qué altura máxima alcanza?
        \item ¿Cuándo alcanza una rapidez de 8.0 m/s?
    \end{enumerate}

    \item \textbf{(Dificultad: Alta) Encuentro de dos objetos.} \\
    Un automóvil y un camión parten del reposo en el mismo instante, con el automóvil cierta distancia detrás del camión. El camión tiene una aceleración constante de 2.10 m/s$^2$, y el automóvil, de 3.40 m/s$^2$. El automóvil rebasa al camión después de que éste ha recorrido 40.0 m.
    \begin{enumerate}[label=\alph*)]
        \item ¿Cuánto tiempo tarda el automóvil en rebasar al camión?
        \item ¿Qué tan atrás del camión estaba el automóvil inicialmente?
        \item ¿Qué rapidez tienen los vehículos cuando están lado a lado?
    \end{enumerate}

    \item \textbf{(Dificultad: Muy Alta) Caída libre y tiempo de reacción.} \\
    Un estudiante deja caer una regla para medir el tiempo de reacción de su compañero de laboratorio. La regla cae 18.0 cm antes de que el compañero la atrape.
    \begin{enumerate}[label=\alph*)]
        \item ¿Cuál fue el tiempo de reacción del compañero?
        \item Después de que el compañero suelta la regla, ¿cuánto tiempo tardará ésta en alcanzar una rapidez igual a la rapidez de un automóvil que viaja a 55 mi/h?
    \end{enumerate}

\end{enumerate}

\end{document}